% C.3bs. Complete restricted extension; unregistered pending full audits.
\paragraph{C.3bs. Component families when the dual complex has no second
cohomology (complete draft).}
Let $k$ be algebraically closed of characteristic $p>0$ and let
$E=\bigcup_{i=1}^cE_i$ be a projective SNC surface with smooth connected
components. Index each connected double curve $C_e$ and each triple
point separately, with the ordered-component orientations. Write
$\Delta$ for this finite two-dimensional dual complex. Assume
\begin{equation}
 H^2(\Delta,\mathbf Q)=0,\qquad
 \rho:\bigoplus_i\operatorname{Pic}(E_i)_{\mathbf Q}
       \longrightarrow\bigoplus_e\operatorname{Pic}(C_e)_{\mathbf Q}
       \quad\text{is surjective},                         \tag{DC.1}
\end{equation}
where $\rho$ is the alternating restriction map of (TF.1).
Let $S(E)=F_\gamma^2K_0(E)_{\mathbf Q}$ and let $J_K(E)$ be the
actual fixed smooth-proper-curve endpoint relation group of EP.
Then
\begin{equation}
 \varphi:S(E)\xrightarrow{\ \sim\ }
           \bigoplus_i\operatorname{CH}_0(E_i)_{\mathbf Q},
 \qquad
 \varphi(J_K(E))=
           \bigoplus_i\operatorname{CH}_0(E_i)^0_{\mathbf Q}.
                                                               \tag{DC.2}
\end{equation}
In particular $S(E)/J_K(E)\simeq\mathbf Q^c$.

Retain the exact normal strict-local threefold scope of C.3bm:
$R$ is the strict henselization at a closed $k$-point of a normal
local ring essentially of finite type over $k$, $\dim R=3$, the
puncture is regular, and $Z\to\operatorname{Spec}R$ is a projective
regular modification which is an isomorphism over that puncture.
Suppose that $E$ is its full reduced SNC closed fibre and that
$D=\sum_i a_iE_i$, with $a_i>0$, satisfies $\mathcal O_E(-D)$ ample.
If every double curve is rational and $H^2(\Delta,\mathbf Q)=0$, then
\begin{equation}
                        B^2(Z\setminus E)_{\mathbf Q}=0. \tag{DC.3}
\end{equation}
This includes triple points and imposes no hypothesis on the
Albanese kernels of the components. It asserts no conclusion when
the second dual cohomology is nonzero, and no higher-dimensional
analogue or new resolution-existence theorem.

\begin{proof}
\emph{1. Keep the full canonical homotopy totalization.}
For a fixed smooth projective parameter curve $P$, or for
$P=\operatorname{Spec}k$, put
\[
 \mathcal C_P^0=\bigoplus_iK(E_i\times P)_{\mathbf Q},\quad
 \mathcal C_P^1=\bigoplus_eK(C_e\times P)_{\mathbf Q},\quad
 \mathcal C_P^2=\bigoplus_vK(P)_{\mathbf Q}.
\]
As in (ET.6), actual closed-union descent gives
\begin{equation}
 \mathcal T_P=\operatorname{fib}(\mathcal C_P^1\to\mathcal C_P^2),
 \qquad
 K(E\times P)_{\mathbf Q}\simeq
                \operatorname{fib}(\mathcal C_P^0\to\mathcal T_P).
                                                               \tag{DC.4}
\end{equation}
The latter map uses the canonical nullhomotopy supplied by repeated
geometric pullback. This is the actual diagram, not the totalization
of its homotopy groups with an arbitrarily chosen zero homotopy.
The exact inputs are the every-degree closed-union excision of
\citet[Theorems A and C, printed pp.~29--30]
{rootGeisserHesselholt2011}, Zariski descent
\citep[Theorem 8.1 and proof, p.~367]{D_ThomasonTrobaugh1990},
and finite induction, as proved in C.3bq.

Because there is no fourth stratum, $H^2(\Delta,\mathbf Q)=0$ says
that the rational matrix
\[
             \delta^1:C^1(\Delta,\mathbf Q)
                              \longrightarrow C^2(\Delta,\mathbf Q)
\]
is surjective. Fix a rational linear right inverse to this matrix.
For every integer $j$, pullback from $P$ gives maps
$K_j(P)_{\mathbf Q}\to K_j(C_e\times P)_{\mathbf Q}$.
Restriction to every marked point of the connected curve $C_e$
is the identity on $K_j(P)_{\mathbf Q}$. Therefore the fixed right
inverse, followed by these constant pullbacks, proves
\begin{equation}
 \pi_j\mathcal C_P^1\longrightarrow\pi_j\mathcal C_P^2
          \quad\text{surjective for every }j.             \tag{DC.5}
\end{equation}
This argument is natural in the parameter, including both endpoint
sections; it makes no assertion about a splitting of the component
map in (DC.4). The long exact sequence of its first fibre now gives
\begin{equation}
 \pi_0\mathcal T_P\hookrightarrow\pi_0\mathcal C_P^1,
 \qquad
 \pi_1\mathcal T_P\hookrightarrow\pi_1\mathcal C_P^1.    \tag{DC.6}
\end{equation}
These injections are what removes the possible coherent vertex
obstruction. Merely arranging supports away from triple points
would not prove them.

\emph{2. The component map is an isomorphism.}
A tuple of pure weight-two component classes restricts to zero in
$K_0$ of every double curve. Its image in $\pi_0\mathcal T_k$
therefore vanishes by (DC.6). Exactness of (DC.4), followed by the
weight-two projector in $K_0(E)_{\mathbf Q}$, lifts the tuple.
This proves surjectivity of $\varphi$.
The projectors exist by the finite gamma filtration on the
projective surface and on its product with a curve
\citep[Chapter II, Proposition 4.9, Theorem 4.10,
and Remark 8.8.5]{D_Weibel2013KbookII}.

If $s\in\ker\varphi$, exactness lifts it to
$\alpha\in\pi_1\mathcal T_k$.
Replace $\alpha$ by $(\psi^2-2)\alpha/2$; its boundary remains $s$.
Adams commutation with these singular-scheme boundaries uses exactly
\citet[p.~258, paragraph preceding Lemma 2.1]{frC3bcLevine1994},
the arbitrary-scheme iterated-fibre and delooping construction,
not the regular-ambient assertion of that lemma.
On each proper smooth double curve the image of the new $\alpha$
is generated by $u([L]-1)$, $u\in k^*$ and $L\in\operatorname{Pic}(C_e)$,
by the full localization and Picard-times-units argument of C.3br
Step~3 \citep[Section 7, Propositions 2.10 and 3.2]{rootQuillen1973}
\citep[Lemma 5.2 and proof, pp.~395--396]{D_PedriniWeibel1994}.
Use surjectivity of $\rho$, one fixed unit at a time, to obtain
$\gamma\in\pi_1\mathcal C_k^0$ with those exact edge restrictions.
The image of $\gamma$ in $\pi_1\mathcal T_k$ equals $\alpha$ by
the second injection in (DC.6). Its boundary is zero, proving $s=0$.
No semisimplicity of $\pi_1\mathcal T_k$ is used.

\emph{3. Actual component families with all edge restrictions corrected.}
Move a component zero-cycle away from its double-curve boundary by
the curve and Riemann--Roch argument in C.3br Step~4. It suffices to
realize $[y]-[x]$ with $x,y\in E_i$ off that boundary.
Choose a high smooth connected projective curve $P\subset E_i$
through $x,y$, transverse to every double curve and avoiding all
triple points. The prescribed-value/first-jet dimension argument
of C.3br applies unchanged; avoiding finitely many additional
points is a simultaneous nonempty open condition
\citep[Lemma 33.47.3 and proof]{D_StacksBertiniUniformHigher}.

Let $o,t\in P(k)$ be the chosen parameter points. The graph perfect
class on $E_i\times P$, projected to weight two and extended as a
tuple with zero on other components, gives $\alpha_P$.
Its alternating edge restriction is the finite sum of transverse
classes $\pm[\mathcal O_c]\boxtimes[\mathcal O_q]$, where every
parameter $q$ differs from $o,t$ and each intersection point $c$
is away from the triple vertices. For each such $q$, lift its whole
edge Picard tuple through $\rho$ and take the pure weight-one
component divisor classes $\ell_{j,q}$. The actual correction
\[
                  v_j=\sum_q\ell_{j,q}\boxtimes[\mathcal O_q]
\]
has zero values at both endpoints and exactly the same edge
restriction as $\alpha_P$.
Thus $\alpha_P-v$ maps to zero in $\pi_0\mathcal C_P^1$.
Its canonical image in $\pi_0\mathcal T_P$ is zero by the first
injection in (DC.6). Exactness of the second fibre in (DC.4)
therefore gives an actual class on the entire fixed $E\times P$.
Its weight-two projection preserves all component restrictions.
The endpoint difference maps to $[y]-[x]$ on the selected component
and zero on all others.

All divisor corrections use finitely many pointwise rational
Picard lifts. A lift of a whole Picard family through $\rho$ is
unnecessary. Canonical vertex coherence is addressed by (DC.5)--(DC.6),
not by treating a zero edge $K_0$ class as a chosen zero complex.

\emph{4. The exact relation group and the puncture consequence.}
These differences generate every component degree-zero Chow group.
Conversely component degree is constant in any actual proper perfect
family, by $\operatorname{ch}_2$ and proper pushforward to the
connected parameter curve. This proves (DC.2).

In the resolved setting, the weighted negative-energy argument of
C.3bq, (ET.9), proves surjectivity of the full joint restriction-degree
map, including all triple intersections. Explicitly an annihilating
functional with normalized edge coefficients $u_{ij}$ satisfies
\[
 0=-\sum_e a_i a_j\deg((-D)|_{C_e})u_{ij}^2
   -\sum_{v=(i,j,k)}a_i a_j a_k
                       (u_{ij}+u_{jk}+u_{ki})^2.
\]
Every coefficient in the first sum is positive, so all edge
coefficients vanish. If the double curves are rational, their
Picard groups tensored with $\mathbf Q$ are exactly their degree
coordinates. Thus this numerical surjectivity is the full $\rho$
surjectivity needed in (DC.1).

Finally C.3bm identifies the ordinary puncture quotient with
$S(E)/(J_K(E)+\theta\operatorname{CH}_1(E)_{\mathbf Q})$.
By (DC.2) only the component degrees survive. The symmetric
exceptional degree matrix of C.3br Step~7 has nonnegative
off-diagonal entries and strictly negative positive-weighted row
sums, so its displayed weighted-square identity makes it invertible.
Exceptional curve classes kill every remaining degree. This proves
(DC.3), with arbitrary component Albanese kernels retained throughout.
\end{proof}

\emph{Scope of the extension.}
The condition $H^2(\Delta,\mathbf Q)=0$ supplies surjectivity on
the whole $K_j(P)$ coefficient group for every fixed parameter,
which is stronger than pointwise vanishing of a period.
When it fails, rational edge Picard groups alone do not supply
the injections (DC.6). The possible fixed-parameter $K_1(P)$
vertex obstruction is not discarded in this statement.

